\section{Introduction}
\label{sec:intro}

If a model's answer to a multiple-choice question depends on where the correct option is printed, then its benchmark score partly measures how the benchmark's answer key lines up with the model's habits.
Multiple-choice accuracy is the default currency for comparing large language models (LLMs), and leaderboards routinely separate models by a few points.
Benchmark builders and anyone comparing models with small accuracy gaps therefore need to know how large the positional component of a score is.

A common form of this worry is directional: models prefer early positions (A/B), so scores are inflated.
We started from a concrete version of this hypothesis, fixed before running any model: models over-select A/B; this inflates accuracy by 3--8 points on 4-option benchmarks; the effect is smallest for true/false ($<$2 points) and largest for 5-option items ($>$10 points); it is strongest on medium-difficulty questions; and near a model's decision boundary up to 15\% of apparent accuracy is positional.

Prior work establishes that LLMs are sensitive to option order and that this sensitivity is tied to label tokens as well as positions~\citep{zheng2024selectors,pezeshkpour2024order,robinson2023mcqa}.
Three things are harder to read off from it.
First, comparing formats across different benchmarks confounds the number of options with question content.
Second, stratifying by difficulty is circular if it reuses the same responses that measure the bias.
Third, ``inflates scores'' is ambiguous: a score can be compared with a balanced key, with a key placed where the model likes to answer, or with the benchmark as published, and these give different numbers.

{\bf How do we separate these?}
We use a within-question design (\secref{sec:method}).
We show each question once per cyclic rotation of its options, so the correct answer occupies every position exactly once.
We build 2-, 3-, 4-, and 5-option versions from the \emph{same} questions of \mmlu~\citep{hendrycks2021mmlu} and \csqa~\citep{talmor2019commonsenseqa}, add native true/false items from \boolq~\citep{clark2019boolq}, and swap A--D labels for 1--4.
We measure difficulty for each model from the \emph{other} models' responses, and we report three explicit definitions of inflation.
We run seven open-weight instruct models from three families (1.5B--14B parameters) and score 298{,}158 prompts.

{\bf What do we find?}
Mostly, the hypothesis does not hold.
Position matters: six of seven models differ by 4--8 points between their best and worst key position on \mmlu.
But the direction differs from model to model, and five of seven models score below their own average when the key is at A.
Scores on the benchmarks as published are 1.3 points (\mmlu) and 0.3 points (\csqa) above position-balanced accuracy, and on \mmlu none of the 1.3 points comes from where the keys are placed.
What position preferences do produce is instability: about 30\% of \mmlu questions and 28\% of \csqa questions are answered correctly under some option orders and incorrectly under others.

In summary, our contributions are:
\begin{itemize}[leftmargin=*,itemsep=0pt,topsep=0pt]
    \item We design a within-question protocol that balances key position by rotation and varies option count and label set on identical questions, and we apply it to seven open-weight models on three benchmarks.
    \item We show that position effects are real and stable within a model but not directed toward A/B, and that as-published scores exceed position-balanced accuracy by about one point or less.
    \item We show that the effect tracks uncertainty: it is near zero on easy questions and peaks where a model is near 50\% accuracy, with a best-versus-worst-position gap of up to 19.7 points.
    \item We find that format matters in unexpected ways: word-answer true/false gives the largest order effect (14.2 points on average), and numeric labels more than double the spread relative to letters (5.9 to 13.6 points).
\end{itemize}

Our results come from open-weight models of at most 14B parameters, five of them from one family; we could not run hosted frontier models, and we discuss this and other limits in \secref{sec:discussion}.
We review related work in \secref{sec:related}, describe the protocol in \secref{sec:method}, present results in \secref{sec:results}, and conclude in \secref{sec:conclusion}.
